\begingroup\setbox0=\vbox{\section*{Trapezia and Their Area}}\endgroup
\chapterhead{13}{Trapezia and Their Area}

\shead{Two parallel bases}
\begin{knowbox}
\begin{minipage}[t]{0.6\linewidth}
A \textbf{trapezium} is a quadrilateral with \emph{one pair of parallel opposite sides}: the \textbf{longer base} and the \textbf{shorter base}. The other two sides are the \textbf{legs}. A trapezium with one leg perpendicular to the bases is a \textbf{right trapezium}.\par
\smallskip
The \textbf{height} is the distance between the bases.
\[S=\frac{(a+b)\times h}{2}\qquad h=\frac{S\times2}{a+b}\qquad a+b=\frac{S\times2}{h}\]
\textbf{The diagonals.} In trapezium ABCD (AB parallel to CD), triangles ACD and BCD share base CD and have equal heights, so $S_{ACD}=S_{BCD}$; removing the common part gives $S_{AOD}=S_{BOC}$.
\end{minipage}\hfill
\begin{minipage}[t]{0.36\linewidth}\vspace{4pt}\centering
\begin{tikzpicture}[scale=0.85]
\coordinate (D) at (0,0); \coordinate (C) at (4.2,0); \coordinate (A) at (0.9,2.2); \coordinate (B) at (2.9,2.2); \coordinate (H) at (0.9,0);
\coordinate (O) at (intersection of A--C and B--D);
\fill[tmBlue!18] (A)--(O)--(D)--cycle; \fill[tmBlue!18] (B)--(O)--(C)--cycle;
\draw[geo] (A)--(B)--(C)--(D)--cycle; \draw[geo] (A)--(C) (B)--(D);
\draw[dashed,tmOrange!80!black,line width=0.7pt] (A)--(H); \ramark{H}{C}{A}
\node[lbl,above] at (A){A}; \node[lbl,above] at (B){B}; \node[lbl,below right] at (C){C}; \node[lbl,below left] at (D){D}; \node[lbl,below=2pt] at (O){O};
\node[dim,left] at (0.9,1.1){$h$}; \node[dim,above=7pt] at (1.9,2.2){$b$}; \node[dim,below=7pt] at (2.6,0){$a$};
\end{tikzpicture}
\end{minipage}
\end{knowbox}
\gap
\begin{exbox}{\ExWord\ 1}
Find the area of a trapezium with bases $12$ cm and $8$ cm and height $5$ cm.
\guide
$S=(12+8)\times5:2=50$ cm$^2$.
\dapso{$50$ cm$^2$.}
\end{exbox}
\gap
\begin{exbox}{\ExWord\ 2}
A trapezium-shaped field has a longer base of $120$ m; the shorter base is $\dfrac23$ of the longer, and the height is the mean of the two bases. Every $100$ m$^2$ yields $64$ kg of paddy. How much paddy does the field yield?
\guide
Shorter base: $120\times\dfrac23=80$ m.\par Height: $(120+80):2=100$ m.\par
Area: $(120+80)\times100:2=10\,000$ m$^2$.\par Paddy: $64\times(10\,000:100)=6\,400$ kg.
\dapso{$6\,400$ kg.}
\end{exbox}
\gap
\begin{exbox}{\ExWord\ 3}
A trapezium has area $60$ cm$^2$ and height $6$ cm; the longer base is $4$ cm longer than the shorter. Find both bases.
\guide
Sum of the bases: $60\times2:6=20$ cm.\par This is a sum-and-difference problem: longer base $(20+4):2=12$ cm, shorter $8$ cm.
\dapso{$12$ cm and $8$ cm.}
\end{exbox}
\gap
\begin{exbox}{\ExWord\ 4}
Trapezium ABCD has shorter base $AB=6$ cm, longer base $CD=9$ cm and height $4$ cm. Find the areas of triangles ABD and BCD, and show that $S_{AOD}=S_{BOC}$ (O is where the diagonals cross).
\guide
Both triangles have the trapezium's height: $S_{ABD}=6\times4:2=12$ cm$^2$ and $S_{BCD}=9\times4:2=18$ cm$^2$.\par
$S_{ACD}=S_{BCD}=18$ cm$^2$ (same base CD, same height).\par Removing their common part, triangle OCD, gives $S_{AOD}=S_{BOC}$.
\dapso{$12$ cm$^2$; $18$ cm$^2$.}
\end{exbox}
\dividerbar{Practice}
\begin{probbox}{1}{\relax}
Find the area of a trapezium with bases $a$, $b$ and height $h$:
\begin{cols2}
\pa{a} $a=14$ cm, $b=6$ cm, $h=7$ cm; & \pa{b} $a=2.8$ m, $b=1.6$ m, $h=0.5$ m;\\
\pa{c} $a=\dfrac34$ dm, $b=\dfrac12$ dm, $h=\dfrac23$ dm; & \pa{d} $a=1.5$ dm, $b=9$ cm, $h=8$ cm.
\end{cols2}
\end{probbox}
\probsol{
\pa{a} $(14+6)\times7:2=70$ cm$^2$.\par \pa{b} $(2.8+1.6)\times0.5:2=1.1$ m$^2$.\par \pa{c} $\left(\dfrac34+\dfrac12\right)\times\dfrac23:2=\dfrac5{12}$ dm$^2$.\par \pa{d} $1.5$ dm $=15$ cm; $(15+9)\times8:2=96$ cm$^2$.
\dapso{$70$ cm$^2$; $1.1$ m$^2$; $\dfrac5{12}$ dm$^2$; $96$ cm$^2$.}}
\gap
\begin{probbox}{2}{\relax}
A right trapezium has bases $10$ cm and $7$ cm, and the leg perpendicular to them is $6$ cm. Find its area.
\end{probbox}
\probsol{
The perpendicular leg is the height: $S=(10+7)\times6:2=51$ cm$^2$.
\dapso{$51$ cm$^2$.}}
\gap
\begin{probbox}{3}{\relax}
A trapezium has area $72$ cm$^2$ and bases $10$ cm and $14$ cm. Find its height.
\end{probbox}
\probsol{
$h=72\times2:(10+14)=6$ cm.
\dapso{$6$ cm.}}
\gap
\begin{probbox}{4}{\relax}
A trapezium has area $45$ m$^2$, height $5$ m and shorter base $7$ m. Find the longer base.
\end{probbox}
\probsol{
Sum of bases: $45\times2:5=18$ m.\par Longer base: $18-7=11$ m.
\dapso{$11$ m.}}
\gap
\begin{probbox}{5}{\relax}
A trapezium-shaped field has bases $60$ m and $40$ m and height $30$ m. Every $100$ m$^2$ yields $70$ kg of paddy. How much paddy does the field yield?
\end{probbox}
\probsol{
Area: $(60+40)\times30:2=1\,500$ m$^2$.\par Paddy: $70\times15=1\,050$ kg.
\dapso{$1\,050$ kg.}}
\gap
\begin{probbox}{6}{\relax}
A trapezium's bases add up to $48$ cm; the shorter is $\dfrac35$ of the longer, and the height is $\dfrac14$ of the longer base. Find its area.
\end{probbox}
\probsol{
Longer base: $48:(3+5)\times5=30$ cm; height: $30\times\dfrac14=7.5$ cm.\par $S=48\times7.5:2=180$ cm$^2$.
\dapso{$180$ cm$^2$.}}
\gap
\begin{probbox}{7}{\relax}
A trapezium has bases $25$ cm and $15$ cm. Lengthening the longer base by $5$ cm adds $20$ cm$^2$ to the area. Find the original area.
\end{probbox}
\probsol{
The added part is a triangle with base $5$ cm and the trapezium's height: $h=20\times2:5=8$ cm.\par $S=(25+15)\times8:2=160$ cm$^2$.
\dapso{$160$ cm$^2$.}}
\gap
\begin{probbox}{8}{\relax}
Rectangle ABCD has $AB=20$ cm and $BC=12$ cm. M is on AB with $AM=8$ cm. Find the area of trapezium MBCD.
\end{probbox}
\probsol{
MBCD has bases $MB=12$ cm and $CD=20$ cm and height $BC=12$ cm: $S=(12+20)\times12:2=192$ cm$^2$.\par
\emph{Another way:} $20\times12-8\times12:2=240-48=192$ cm$^2$.
\dapso{$192$ cm$^2$.}}
\gap
\begin{probbox}{9}{\relax}
Trapezium ABCD has $AB=4$ cm, $CD=10$ cm and height $6$ cm. Find the areas of triangles ABC and ACD.
\end{probbox}
\probsol{
$S_{ABC}=4\times6:2=12$ cm$^2$;\par $S_{ACD}=10\times6:2=30$ cm$^2$.\par \emph{Check:} $12+30=42=(4+10)\times6:2$.
\dapso{$12$ cm$^2$; $30$ cm$^2$.}}
\gap
\begin{probbox}{10}{\relax}
Find the area of a trapezium with bases $15$ cm and $9$ cm and height $6$ cm.
\end{probbox}
\probsol{
$S=(15+9)\times6:2=72$ cm$^2$.
\dapso{$72$ cm$^2$.}}
\gap
\begin{probbox}{11}{\relax}
Find the area of a trapezium with bases $6.5$ cm and $3.5$ cm and height $4.2$ cm.
\end{probbox}
\probsol{
$S=(6.5+3.5)\times4.2:2=21$ cm$^2$.
\dapso{$21$ cm$^2$.}}
\gap
\begin{probbox}{12}{\relax}
Find the area of a trapezium with bases $\dfrac34$ m and $\dfrac12$ m and height $\dfrac23$ m.
\end{probbox}
\probsol{
$S=\left(\dfrac34+\dfrac12\right)\times\dfrac23:2=\dfrac54\times\dfrac23:2=\dfrac5{12}$ m$^2$.
\dapso{$\dfrac5{12}$ m$^2$.}}
\gap
\begin{probbox}{13}{\relax}
Find the area of a trapezium with bases $14$ dm and $10$ dm and height $7.5$ dm.
\end{probbox}
\probsol{
$S=(14+10)\times7.5:2=90$ dm$^2$.
\dapso{$90$ dm$^2$.}}
\gap
\begin{probbox}{14}{\relax}
Find the area of a trapezium with bases $1.2$ m and $80$ cm and height $5$ dm.
\end{probbox}
\probsol{
In decimetres: $1.2$ m $=12$ dm; $80$ cm $=8$ dm.\par $S=(12+8)\times5:2=50$ dm$^2$.
\dapso{$50$ dm$^2$.}}
\gap
\begin{probbox}{15}{\relax}
A trapezium has area $84$ cm$^2$ and bases $16$ cm and $12$ cm. Find its height.
\end{probbox}
\probsol{
Height $=$ area $\times2\,:$ sum of bases.\par $h=84\times2:(16+12)=6$ cm.
\dapso{$6$ cm.}}
\gap
\begin{probbox}{16}{\relax}
A trapezium has area $45$ cm$^2$ and bases $8.5$ cm and $6.5$ cm. Find its height.
\end{probbox}
\probsol{
Height $=$ area $\times2\,:$ sum of bases.\par $h=45\times2:(8.5+6.5)=6$ cm.
\dapso{$6$ cm.}}
\gap
\begin{probbox}{17}{\relax}
A trapezium has area $150$ cm$^2$ and bases $23$ cm and $17$ cm. Find its height.
\end{probbox}
\probsol{
Height $=$ area $\times2\,:$ sum of bases.\par $h=150\times2:(23+17)=7.5$ cm.
\dapso{$7.5$ cm.}}
\gap
\begin{probbox}{18}{\relax}
A trapezium has area $96$ cm$^2$, height $8$ cm and shorter base $10$ cm. Find the longer base.
\end{probbox}
\probsol{
Sum of bases: $96\times2:8=24$ cm.\par Longer base: $24-10=14$ cm.
\dapso{$14$ cm.}}
\gap
\begin{probbox}{19}{\relax}
A trapezium has area $63$ cm$^2$, height $7$ cm and shorter base $7.5$ cm. Find the longer base.
\end{probbox}
\probsol{
Sum of bases: $63\times2:7=18$ cm.\par Longer base: $18-7.5=10.5$ cm.
\dapso{$10.5$ cm.}}
\gap
\begin{probbox}{20}{\relax}
A trapezium has area $110$ cm$^2$, height $10$ cm and shorter base $9$ cm. Find the longer base.
\end{probbox}
\probsol{
Sum of bases: $110\times2:10=22$ cm.\par Longer base: $22-9=13$ cm.
\dapso{$13$ cm.}}
\gap
\begin{probbox}{21}{\relax}
A trapezium-shaped field has bases $80$ m and $50$ m and height $40$ m. Every $100$ m$^2$ yields $65$ kg of paddy. What is the total yield?
\end{probbox}
\probsol{
Area: $(80+50)\times40:2=2\,600$ m$^2$.\par Yield: $2\,600:100\times65=1\,690$ kg.
\dapso{$1\,690$ kg.}}
\gap
\begin{probbox}{22}{\relax}
A trapezium-shaped field has bases $150$ m and $90$ m and height $60$ m. Every $100$ m$^2$ yields $60$ kg of paddy. What is the total yield?
\end{probbox}
\probsol{
Area: $(150+90)\times60:2=7\,200$ m$^2$.\par Yield: $7\,200:100\times60=4\,320$ kg.
\dapso{$4\,320$ kg.}}
\gap
\begin{probbox}{23}{\relax}
A trapezium's bases add up to $36$ cm and differ by $8$ cm; its height is $10$ cm. Find each base and the area.
\end{probbox}
\probsol{
Longer base: $(36+8):2=22$ cm; shorter: $14$ cm.\par Area: $36\times10:2=180$ cm$^2$.
\dapso{Bases $22$ cm and $14$ cm; $180$ cm$^2$.}}
\gap
\begin{probbox}{24}{\relax}
A trapezium's bases add up to $50$ cm and differ by $14$ cm; its height is $12$ cm. Find each base and the area.
\end{probbox}
\probsol{
Longer base: $(50+14):2=32$ cm; shorter: $18$ cm.\par Area: $50\times12:2=300$ cm$^2$.
\dapso{Bases $32$ cm and $18$ cm; $300$ cm$^2$.}}
\gap
\begin{probbox}{25}{\relax}
A trapezium-shaped plot has longer base $60$ m; the shorter base is $\dfrac{2}{3}$ of it and the height is $\dfrac{1}{2}$ of it. Find the area.
\end{probbox}
\probsol{
Shorter base: $60\times\dfrac{2}{3}=40$ m.\par Height: $60\times\dfrac{1}{2}=30$ m.\par Area: $(60+40)\times30:2=1\,500$ m$^2$.
\dapso{$1\,500$ m$^2$.}}
\gap
\begin{probbox}{26}{\relax}
A trapezium-shaped plot has longer base $48$ m; the shorter base is $\dfrac{3}{4}$ of it and the height is $\dfrac{5}{8}$ of it. Find the area.
\end{probbox}
\probsol{
Shorter base: $48\times\dfrac{3}{4}=36$ m.\par Height: $48\times\dfrac{5}{8}=30$ m.\par Area: $(48+36)\times30:2=1\,260$ m$^2$.
\dapso{$1\,260$ m$^2$.}}
\gap
\begin{probbox}{27}{\relax}
Trapezium ABCD has right angles at A and D, $AB=9$ cm, $CD=15$ cm and $AD=8$ cm. Find its area.
\end{probbox}
\probsol{
AD is perpendicular to both bases, so it is the height.\par $S=(9+15)\times8:2=96$ cm$^2$.
\dapso{$96$ cm$^2$.}}
\gap
\begin{probbox}{28}{\relax}
Trapezium ABCD has right angles at A and D, $AB=4.5$ cm, $CD=7.5$ cm and $AD=6$ cm. Find its area.
\end{probbox}
\probsol{
AD is perpendicular to both bases, so it is the height.\par $S=(4.5+7.5)\times6:2=36$ cm$^2$.
\dapso{$36$ cm$^2$.}}
\gap
\begin{probbox}{29}{\relax}
Rectangle ABCD has $AB=25$ cm and $BC=14$ cm. M on AB has $AM=9$ cm. Find the areas of trapezium MBCD and triangle AMD.
\end{probbox}
\probsol{
MBCD has bases $MB=16$ cm and $CD=25$ cm and height $BC=14$ cm.\par $S_{MBCD}=(16+25)\times14:2=287$ cm$^2$.\par $S_{AMD}=9\times14:2=63$ cm$^2$.
\dapso{$S_{MBCD}=287$ cm$^2$, $S_{AMD}=63$ cm$^2$.}}
\gap
\begin{probbox}{30}{\relax}
Rectangle ABCD has $AB=30$ cm and $BC=16$ cm. M on AB has $AM=12$ cm. Find the areas of trapezium MBCD and triangle AMD.
\end{probbox}
\probsol{
MBCD has bases $MB=18$ cm and $CD=30$ cm and height $BC=16$ cm.\par $S_{MBCD}=(18+30)\times16:2=384$ cm$^2$.\par $S_{AMD}=12\times16:2=96$ cm$^2$.
\dapso{$S_{MBCD}=384$ cm$^2$, $S_{AMD}=96$ cm$^2$.}}
\dividerbar{Challenge -- Extension}
\begin{probbox}{31}{\stars{1}}
A trapezium has bases $12$ cm and $18$ cm. Lengthening the shorter base by $6$ cm turns it into a rectangle and adds $27$ cm$^2$. Find the trapezium's area.
\end{probbox}
\probsol{
The added part is a triangle with side $6$ cm and the trapezium's height: $h=27\times2:6=9$ cm.\par $S=(12+18)\times9:2=135$ cm$^2$.
\dapso{$135$ cm$^2$.}}
\gap
\begin{probbox}{32}{\stars{1}}
A trapezium-shaped field has bases adding up to $140$ m, the longer $20$ m longer than the shorter. Lengthening each base by $10$ m adds $500$ m$^2$. Find the original area and both bases.
\end{probbox}
\probsol{
The bases grow by $20$ m in total, so $h=500\times2:20=50$ m.\par $S=140\times50:2=3\,500$ m$^2$.\par
Longer base: $(140+20):2=80$ m; shorter: $60$ m.
\dapso{$3\,500$ m$^2$; $80$ m and $60$ m.}}
\gap
\begin{probbox}{33}{\stars{1}}
In trapezium ABCD, the shorter base AB is $\dfrac13$ of the longer base CD, and $S_{ABD}=15$ cm$^2$. Find the area of the trapezium.
\end{probbox}
\probsol{
ABD and BCD both have the trapezium's height, and $CD=3\times AB$, so $S_{BCD}=15\times3=45$ cm$^2$.\par $S_{ABCD}=15+45=60$ cm$^2$.
\dapso{$60$ cm$^2$.}}
\gap
\begin{probbox}{34}{\stars{2}}
Right trapezium ABCD has right angles at A and D, $AB=6$ cm, $CD=12$ cm and $AD=8$ cm. M is the midpoint of AD. Find $S_{MBC}$.
\end{probbox}
\probsol{
$S_{ABCD}=(6+12)\times8:2=72$;\par $S_{ABM}=6\times4:2=12$;\par $S_{MDC}=12\times4:2=24$ cm$^2$.\par
$S_{MBC}=72-12-24=36$ cm$^2$, exactly half the trapezium.
\dapso{$36$ cm$^2$.}}
\gap
\begin{probbox}{35}{\stars{2}}
In trapezium ABCD (AB parallel to CD), the diagonals cross at O. $S_{AOB}=4$ cm$^2$ and $S_{BOC}=8$ cm$^2$. Find the area of the trapezium.
\end{probbox}
\probsol{
AOB and BOC share the height from B, so $OC=2\times OA$.\par
By the trapezium property, $S_{AOD}=S_{BOC}=8$ cm$^2$.\par COD and AOD share the height from D and $OC=2\times OA$, so $S_{COD}=16$ cm$^2$.\par
$S_{ABCD}=4+8+8+16=36$ cm$^2$.
\dapso{$36$ cm$^2$.}}
\gap
\begin{probbox}{36}{\stars{1}}
A trapezium has bases $30$ cm and $18$ cm. Lengthening the longer base by $6$ cm (same height) adds $27$ cm$^2$. Find the original area.
\end{probbox}
\probsol{
The extra part is a triangle with base $6$ cm and the trapezium's height.\par Height: $27\times2:6=9$ cm.\par Area: $(30+18)\times9:2=216$ cm$^2$.
\dapso{$216$ cm$^2$.}}
\gap
\begin{probbox}{37}{\stars{1}}
A trapezium has bases $22$ cm and $14$ cm. Lengthening the longer base by $4$ cm (same height) adds $18$ cm$^2$. Find the original area.
\end{probbox}
\probsol{
The extra part is a triangle with base $4$ cm and the trapezium's height.\par Height: $18\times2:4=9$ cm.\par Area: $(22+14)\times9:2=162$ cm$^2$.
\dapso{$162$ cm$^2$.}}
\gap
\begin{probbox}{38}{\stars{1}}
A trapezium-shaped garden has bases summing to $90$ m, differing by $30$ m. Adding $5$ m to each base (same height) adds $200$ m$^2$. Find the original area.
\end{probbox}
\probsol{
The extra part is a trapezium with bases $5$ m and $5$ m and the same height.\par Height: $200\times2:(5+5)=40$ m.\par Original area: $90\times40:2=1\,800$ m$^2$.
\dapso{$1\,800$ m$^2$.}}
\gap
\begin{probbox}{39}{\stars{1}}
A trapezium-shaped garden has bases summing to $120$ m, differing by $40$ m. Adding $8$ m to each base (same height) adds $320$ m$^2$. Find the original area.
\end{probbox}
\probsol{
The extra part is a trapezium with bases $8$ m and $8$ m and the same height.\par Height: $320\times2:(8+8)=40$ m.\par Original area: $120\times40:2=2\,400$ m$^2$.
\dapso{$2\,400$ m$^2$.}}
\gap
\begin{probbox}{40}{\stars{1}}
In trapezium ABCD the shorter base AB is $\dfrac{1}{2}$ of the longer base CD. Triangle ABD has area $12$ cm$^2$. Find the area of the trapezium.
\end{probbox}
\probsol{
ABD and BCD have the same height (the trapezium's height).\par So $S_{ABD}=\dfrac{1}{2}S_{BCD}$.\par $S_{BCD}=12:\dfrac{1}{2}=24$ cm$^2$.\par $S_{ABCD}=36$ cm$^2$.
\dapso{$36$ cm$^2$.}}
\gap
\begin{probbox}{41}{\stars{1}}
In trapezium ABCD the shorter base AB is $\dfrac{2}{3}$ of the longer base CD. Triangle ABD has area $18$ cm$^2$. Find the area of the trapezium.
\end{probbox}
\probsol{
ABD and BCD have the same height (the trapezium's height).\par So $S_{ABD}=\dfrac{2}{3}S_{BCD}$.\par $S_{BCD}=18:\dfrac{2}{3}=27$ cm$^2$.\par $S_{ABCD}=45$ cm$^2$.
\dapso{$45$ cm$^2$.}}
\gap
\begin{probbox}{42}{\stars{2}}
Trapezium ABCD (AB parallel to CD) has diagonals meeting at O. $S_{AOB}=4$ cm$^2$ and $S_{BOC}=6$ cm$^2$. Find the area of the trapezium.
\end{probbox}
\probsol{
$S_{ABD}=S_{ABC}$ (same base AB, same height); removing AOB gives $S_{AOD}=S_{BOC}=6$ cm$^2$.\par AOB and BOC share the height from B, so $\dfrac{OA}{OC}=\dfrac{4}{6}$.\par AOD and COD share the height from D, so $S_{COD}=6\times\dfrac{6}{4}=9$ cm$^2$.\par $S_{ABCD}=4+6+6+9=25$ cm$^2$.
\dapso{$25$ cm$^2$.}}
\gap
\begin{probbox}{43}{\stars{2}}
Trapezium ABCD (AB parallel to CD) has diagonals meeting at O. $S_{AOB}=9$ cm$^2$ and $S_{BOC}=12$ cm$^2$. Find the area of the trapezium.
\end{probbox}
\probsol{
$S_{ABD}=S_{ABC}$ (same base AB, same height); removing AOB gives $S_{AOD}=S_{BOC}=12$ cm$^2$.\par AOB and BOC share the height from B, so $\dfrac{OA}{OC}=\dfrac{9}{12}$.\par AOD and COD share the height from D, so $S_{COD}=12\times\dfrac{12}{9}=16$ cm$^2$.\par $S_{ABCD}=9+12+12+16=49$ cm$^2$.
\dapso{$49$ cm$^2$.}}
\gap
\begin{probbox}{44}{\stars{2}}
Trapezium ABCD (AB parallel to CD) has diagonals meeting at O. $S_{AOB}=2$ cm$^2$ and $S_{BOC}=6$ cm$^2$. Find the area of the trapezium.
\end{probbox}
\probsol{
$S_{ABD}=S_{ABC}$ (same base AB, same height); removing AOB gives $S_{AOD}=S_{BOC}=6$ cm$^2$.\par AOB and BOC share the height from B, so $\dfrac{OA}{OC}=\dfrac{2}{6}$.\par AOD and COD share the height from D, so $S_{COD}=6\times\dfrac{6}{2}=18$ cm$^2$.\par $S_{ABCD}=2+6+6+18=32$ cm$^2$.
\dapso{$32$ cm$^2$.}}
\gap
\begin{probbox}{45}{\stars{1}}
Trapezium ABCD has right angles at A and D, $AB=8$ cm, $CD=14$ cm, $AD=10$ cm. M is the midpoint of AD. Find the area of triangle MBC.
\end{probbox}
\probsol{
$S_{ABCD}=(8+14)\times10:2=110$ cm$^2$.\par $AM=MD=5$ cm.\par $S_{MAB}=20$, $S_{MDC}=35$ cm$^2$.\par $S_{MBC}=110-20-35=55$ cm$^2$.
\dapso{$55$ cm$^2$.}}
\gap
\begin{probbox}{46}{\stars{1}}
Trapezium ABCD has right angles at A and D, $AB=5$ cm, $CD=11$ cm, $AD=12$ cm. M is the midpoint of AD. Find the area of triangle MBC.
\end{probbox}
\probsol{
$S_{ABCD}=(5+11)\times12:2=96$ cm$^2$.\par $AM=MD=6$ cm.\par $S_{MAB}=15$, $S_{MDC}=33$ cm$^2$.\par $S_{MBC}=96-15-33=48$ cm$^2$.
\dapso{$48$ cm$^2$.}}
\gap
\begin{probbox}{47}{\stars{1}}
A right trapezium has bases $10$ cm and $16$ cm. Lengthening the shorter base by $6$ cm makes a rectangle $24$ cm$^2$ larger. Find the trapezium's area.
\end{probbox}
\probsol{
The added piece is a right triangle with legs $6$ cm and the height.\par Height: $24\times2:6=4$ cm.\par Area: $(10+16)\times4:2=52$ cm$^2$.
\dapso{$52$ cm$^2$.}}
\gap
\begin{probbox}{48}{\stars{1}}
A right trapezium has bases $14$ cm and $20$ cm. Lengthening the shorter base by $6$ cm makes a rectangle $27$ cm$^2$ larger. Find the trapezium's area.
\end{probbox}
\probsol{
The added piece is a right triangle with legs $6$ cm and the height.\par Height: $27\times2:6=4.5$ cm.\par Area: $(14+20)\times4.5:2=76.5$ cm$^2$.
\dapso{$76.5$ cm$^2$.}}
\gap
\begin{probbox}{49}{\stars{1}}
A trapezium with bases $20$ cm and $16$ cm has the same area as a square of side $12$ cm. Find its height.
\end{probbox}
\probsol{
Square: $12\times12=144$ cm$^2$.\par Height: $144\times2:(20+16)=8$ cm.
\dapso{$8$ cm.}}
\gap
\begin{probbox}{50}{\stars{1}}
Diagonal AC splits trapezium ABCD (AB parallel to CD) into triangles with $S_{ABC}=24$ cm$^2$ and $S_{ACD}=40$ cm$^2$. The height is $8$ cm. Find both bases.
\end{probbox}
\probsol{
Both triangles have the trapezium's height, $8$ cm.\par $AB=24\times2:8=6$ cm.\par $CD=40\times2:8=10$ cm.
\dapso{$AB=6$ cm, $CD=10$ cm.}}
\clearpage
\titlehead
\begingroup\setbox0=\vbox{\section*{Circles: Circumference and Area}}\endgroup
\chapterhead{14}{Circles: Circumference and Area}

\shead{Centre, radius, diameter}
\begin{knowbox}
\begin{minipage}[t]{0.64\linewidth}
Every point on a \textbf{circle} is the same distance from the \textbf{centre} O. A segment from the centre to the circle is a \textbf{radius} $r$; a segment through the centre joining two points of the circle is a \textbf{diameter} $d$, and $d=r\times2$.\par
\smallskip
We use the number $3.14$ (called ``pi'') to calculate:
\[C=d\times3.14=r\times2\times3.14\qquad\qquad S=r\times r\times3.14\]
\begin{itemize}[leftmargin=1.4em,itemsep=2pt]
\item Given the circumference: $d=C:3.14$.
\item Given the area: find $r\times r=S:3.14$, then the number $r$.
\item If the radius is multiplied by $k$, the circumference is multiplied by $k$ and the area by $k\times k$.
\end{itemize}
\end{minipage}\hfill
\begin{minipage}[t]{0.32\linewidth}\vspace{6pt}\centering
\begin{tikzpicture}[scale=0.85]
\fill[tmBlue!12] (0,0) circle (1.5); \draw[geo] (0,0) circle (1.5);
\draw[line width=0.8pt,tmOrange!80!black] (-1.5,0)--(1.5,0);
\draw[line width=0.8pt,tmGreen!70!black] (0,0)--(60:1.5);
\fill (0,0) circle (1.5pt); \node[lbl,below] at (0,0){O};
\node[dim,above left] at (40:0.7){$r$}; \node[dim,below] at (-0.8,0){$d$};
\end{tikzpicture}
\end{minipage}
\end{knowbox}
\gap
\begin{exbox}{\ExWord\ 1}
Find the circumference and area of a circle of radius $5$ cm.
\guide
$C=5\times2\times3.14=31.4$ cm.\par $S=5\times5\times3.14=78.5$ cm$^2$.
\dapso{$31.4$ cm; $78.5$ cm$^2$.}
\end{exbox}
\gap
\begin{exbox}{\ExWord\ 2}
A circle has circumference $18.84$ dm. Find its area.
\guide
Diameter: $18.84:3.14=6$ dm; radius: $3$ dm.\par $S=3\times3\times3.14=28.26$ dm$^2$.
\dapso{$28.26$ dm$^2$.}
\end{exbox}
\gap
\begin{exbox}{\ExWord\ 3}
\begin{minipage}[t]{0.66\linewidth}
A square has side $10$ cm. A circle inside it touches all four sides. Find the shaded area (the part of the square outside the circle).
\end{minipage}\hfill
\begin{minipage}[t]{0.3\linewidth}\vspace{0pt}\centering
\begin{tikzpicture}[scale=0.22]
\fill[shade] (0,0) rectangle (10,10); \fill[white] (5,5) circle (5);
\draw[geo] (0,0) rectangle (10,10); \draw[geo] (5,5) circle (5);
\draw[|<->|,Slate] (0,-1.2)--node[fill=white,dim]{$10$ cm}(10,-1.2);
\end{tikzpicture}
\end{minipage}
\guide
The circle's diameter equals the side of the square: $d=10$ cm, $r=5$ cm.\par
$S_\text{square}=10\times10=100$ cm$^2$;\par $S_\text{circle}=5\times5\times3.14=78.5$ cm$^2$.\par
Shaded part: $100-78.5=21.5$ cm$^2$.
\dapso{$21.5$ cm$^2$.}
\end{exbox}
\gap
\begin{exbox}{\ExWord\ 4}
A cart wheel has a diameter of $0.6$ m. How far does the cart travel when the wheel turns $100$ times?
\guide
Each turn moves the cart by the wheel's circumference: $0.6\times3.14=1.884$ m.\par
$100$ turns: $1.884\times100=188.4$ m.
\dapso{$188.4$ m.}
\end{exbox}
\dividerbar{Practice}
\begin{probbox}{1}{\relax}
Find the circumference and area of a circle with:
\begin{cols3}
\pa{a} $r=3$ cm; & \pa{b} $r=0.5$ dm; & \pa{c} $d=8$ m.
\end{cols3}
\end{probbox}
\probsol{
\pa{a} $C=18.84$ cm; $S=28.26$ cm$^2$.\par \pa{b} $C=3.14$ dm; $S=0.785$ dm$^2$.\par \pa{c} $C=25.12$ m; $r=4$ m, $S=50.24$ m$^2$.
\dapso{See above.}}
\gap
\begin{probbox}{2}{\relax}
A circle has circumference $12.56$ cm. Find its radius and area.
\end{probbox}
\probsol{
$d=12.56:3.14=4$ cm; $r=2$ cm.\par $S=2\times2\times3.14=12.56$ cm$^2$.
\dapso{$r=2$ cm; $S=12.56$ cm$^2$.}}
\gap
\begin{probbox}{3}{\relax}
A bicycle wheel has a diameter of $65$ cm. How many kilometres does the bicycle travel in $1\,000$ turns of the wheel?
\end{probbox}
\probsol{
Each turn: $65\times3.14=204.1$ cm.\par $1\,000$ turns: $204\,100$ cm $=2.041$ km.
\dapso{$2.041$ km.}}
\gap
\begin{probbox}{4}{\relax}
A car wheel travels $1\,099$ m in $500$ turns. Find its diameter.
\end{probbox}
\probsol{
Circumference: $1\,099:500=2.198$ m.\par Diameter: $2.198:3.14=0.7$ m.
\dapso{$0.7$ m.}}
\gap
\begin{probbox}{5}{\relax}
The minute hand of a clock is $10$ cm long. How far does its tip travel in $1$ hour? In $15$ minutes?
\end{probbox}
\probsol{
In $1$ hour the tip goes once round: $10\times2\times3.14=62.8$ cm.\par In $15$ minutes, a quarter turn: $62.8:4=15.7$ cm.
\dapso{$62.8$ cm; $15.7$ cm.}}
\gap
\begin{probbox}{6}{\relax}
Find the area and perimeter of a semicircle with diameter $8$ cm.
\end{probbox}
\probsol{
$S=4\times4\times3.14:2=25.12$ cm$^2$.\par The perimeter is half the circle plus the diameter: $8\times3.14:2+8=20.56$ cm.
\dapso{$25.12$ cm$^2$; $20.56$ cm.}}
\gap
\begin{probbox}{7}{\relax}
A round flower bed has a radius of $10$ m. A path $1$ m wide is built all round it. Find the area of the path.
\end{probbox}
\probsol{
The outer circle has radius $11$ m.\par Path: $11\times11\times3.14-10\times10\times3.14=(121-100)\times3.14=65.94$ m$^2$.
\dapso{$65.94$ m$^2$.}}
\gap
\begin{probbox}{8}{\relax}
A square sits inside a circle of radius $5$ cm with its four corners on the circle. Find the area of the part of the circle outside the square.
\end{probbox}
\probsol{
The square's diagonals are diameters, $10$ cm long, and split it into $4$ right-angled triangles with sides $5$ cm: $S_\text{square}=4\times5\times5:2=50$ cm$^2$.\par
$S_\text{circle}=78.5$ cm$^2$.\par Required part: $78.5-50=28.5$ cm$^2$.
\dapso{$28.5$ cm$^2$.}}
\gap
\begin{probbox}{9}{\relax}
If the radius of a circle is multiplied by $3$, by how many times are its circumference and area multiplied? Check with $r=1$ cm.
\end{probbox}
\probsol{
The circumference is multiplied by $3$ and the area by $3\times3=9$.\par Example: $r=1$: $C=6.28$, $S=3.14$; $r=3$: $C=18.84=6.28\times3$, $S=28.26=3.14\times9$.
\dapso{$3$ times; $9$ times.}}
\gap
\begin{probbox}{10}{\relax}
Find the circumference and area of a circle of radius $3$ cm.
\end{probbox}
\probsol{
$C=3\times2\times3.14=18.84$ cm.\par $S=3\times3\times3.14=28.26$ cm$^2$.
\dapso{$C=18.84$ cm; $S=28.26$ cm$^2$.}}
\gap
\begin{probbox}{11}{\relax}
Find the circumference and area of a circle of radius $2.5$ cm.
\end{probbox}
\probsol{
$C=2.5\times2\times3.14=15.7$ cm.\par $S=2.5\times2.5\times3.14=19.625$ cm$^2$.
\dapso{$C=15.7$ cm; $S=19.625$ cm$^2$.}}
\gap
\begin{probbox}{12}{\relax}
Find the circumference and area of a circle of radius $0.8$ m.
\end{probbox}
\probsol{
$C=0.8\times2\times3.14=5.024$ m.\par $S=0.8\times0.8\times3.14=2.0096$ m$^2$.
\dapso{$C=5.024$ m; $S=2.0096$ m$^2$.}}
\gap
\begin{probbox}{13}{\relax}
Find the circumference and area of a circle of radius $12$ dm.
\end{probbox}
\probsol{
$C=12\times2\times3.14=75.36$ dm.\par $S=12\times12\times3.14=452.16$ dm$^2$.
\dapso{$C=75.36$ dm; $S=452.16$ dm$^2$.}}
\gap
\begin{probbox}{14}{\relax}
Find the circumference and area of a circle of diameter $10$ cm.
\end{probbox}
\probsol{
$C=10\times3.14=31.4$ cm.\par Radius: $10:2=5$ cm.\par $S=5\times5\times3.14=78.5$ cm$^2$.
\dapso{$C=31.4$ cm; $S=78.5$ cm$^2$.}}
\gap
\begin{probbox}{15}{\relax}
Find the circumference and area of a circle of diameter $4.2$ cm.
\end{probbox}
\probsol{
$C=4.2\times3.14=13.188$ cm.\par Radius: $4.2:2=2.1$ cm.\par $S=2.1\times2.1\times3.14=13.8474$ cm$^2$.
\dapso{$C=13.188$ cm; $S=13.8474$ cm$^2$.}}
\gap
\begin{probbox}{16}{\relax}
Find the circumference and area of a circle of diameter $2.4$ m.
\end{probbox}
\probsol{
$C=2.4\times3.14=7.536$ m.\par Radius: $2.4:2=1.2$ m.\par $S=1.2\times1.2\times3.14=4.5216$ m$^2$.
\dapso{$C=7.536$ m; $S=4.5216$ m$^2$.}}
\gap
\begin{probbox}{17}{\relax}
A circle has circumference $43.96$ cm. Find its radius and area.
\end{probbox}
\probsol{
Diameter: $43.96:3.14=14$ cm.\par Radius: $14:2=7$ cm.\par $S=7\times7\times3.14=153.86$ cm$^2$.
\dapso{$r=7$ cm; $S=153.86$ cm$^2$.}}
\gap
\begin{probbox}{18}{\relax}
A circle has circumference $21.98$ cm. Find its radius and area.
\end{probbox}
\probsol{
Diameter: $21.98:3.14=7$ cm.\par Radius: $7:2=3.5$ cm.\par $S=3.5\times3.5\times3.14=38.465$ cm$^2$.
\dapso{$r=3.5$ cm; $S=38.465$ cm$^2$.}}
\gap
\begin{probbox}{19}{\relax}
A circle has circumference $6.28$ m. Find its radius and area.
\end{probbox}
\probsol{
Diameter: $6.28:3.14=2$ m.\par Radius: $2:2=1$ m.\par $S=1\times1\times3.14=3.14$ m$^2$.
\dapso{$r=1$ m; $S=3.14$ m$^2$.}}
\gap
\begin{probbox}{20}{\relax}
A circle has area $113.04$ cm$^2$. Find its radius and circumference.
\end{probbox}
\probsol{
$r\times r=113.04:3.14=36$.\par Since $6\times6=36$, $r=6$ cm.\par $C=6\times2\times3.14=37.68$ cm.
\dapso{$r=6$ cm; $C=37.68$ cm.}}
\gap
\begin{probbox}{21}{\relax}
A circle has area $12.56$ m$^2$. Find its radius and circumference.
\end{probbox}
\probsol{
$r\times r=12.56:3.14=4$.\par Since $2\times2=4$, $r=2$ m.\par $C=2\times2\times3.14=12.56$ m.
\dapso{$r=2$ m; $C=12.56$ m.}}
\gap
\begin{probbox}{22}{\relax}
A wheel has diameter $0.5$ m. How far does it travel in $200$ turns?
\end{probbox}
\probsol{
One turn: $0.5\times3.14=1.57$ m.\par $200$ turns: $1.57\times200=314$ m.
\dapso{$314$ m.}}
\gap
\begin{probbox}{23}{\relax}
A wheel has diameter $0.7$ m. How far does it travel in $1\,000$ turns?
\end{probbox}
\probsol{
One turn: $0.7\times3.14=2.198$ m.\par $1\,000$ turns: $2.198\times1\,000=2\,198$ m.
\dapso{$2\,198$ m.}}
\gap
\begin{probbox}{24}{\relax}
A bicycle wheel travels $816.4$ m in $400$ turns. Find its diameter.
\end{probbox}
\probsol{
One turn: $816.4:400=2.041$ m.\par Diameter: $2.041:3.14=0.65$ m.
\dapso{$0.65$ m.}}
\gap
\begin{probbox}{25}{\relax}
Find the area and perimeter of a semicircle with diameter $10$ cm.
\end{probbox}
\probsol{
Area: $5\times5\times3.14:2=39.25$ cm$^2$.\par Arc: $10\times3.14:2=15.7$ cm.\par Perimeter (with the diameter): $15.7+10=25.7$ cm.
\dapso{$S=39.25$ cm$^2$; perimeter $25.7$ cm.}}
\gap
\begin{probbox}{26}{\relax}
Find the area and perimeter of a semicircle with diameter $6$ cm.
\end{probbox}
\probsol{
Area: $3\times3\times3.14:2=14.13$ cm$^2$.\par Arc: $6\times3.14:2=9.42$ cm.\par Perimeter (with the diameter): $9.42+6=15.42$ cm.
\dapso{$S=14.13$ cm$^2$; perimeter $15.42$ cm.}}
\gap
\begin{probbox}{27}{\relax}
A ceiling-fan blade is $0.6$ m long from the centre. Find the area of the circle its tip sweeps out.
\end{probbox}
\probsol{
Radius $r=0.6$ m.\par $S=0.6\times0.6\times3.14=1.1304$ m$^2$.
\dapso{$1.1304$ m$^2$.}}
\gap
\begin{probbox}{28}{\relax}
A round pond has diameter $20$ m. Fencing costs $50\,000$ dong per metre. How much does it cost to fence the pond?
\end{probbox}
\probsol{
Circumference: $20\times3.14=62.8$ m.\par Cost: $50\,000\times62.8=3\,140\,000$ dong.
\dapso{$3\,140\,000$ dong.}}
\gap
\begin{probbox}{29}{\relax}
Find the area of a quarter circle of radius $4$ cm.
\end{probbox}
\probsol{
Whole circle: $4\times4\times3.14=50.24$ cm$^2$.\par Quarter: $50.24:4=12.56$ cm$^2$.
\dapso{$12.56$ cm$^2$.}}
\gap
\begin{probbox}{30}{\relax}
Find the area of a quarter circle of radius $10$ cm.
\end{probbox}
\probsol{
Whole circle: $10\times10\times3.14=314$ cm$^2$.\par Quarter: $314:4=78.5$ cm$^2$.
\dapso{$78.5$ cm$^2$.}}
\dividerbar{Challenge -- Extension}
\begin{probbox}{31}{\stars{1}}
A circle has area $50.24$ cm$^2$. Find its circumference.
\end{probbox}
\probsol{
$r\times r=50.24:3.14=16$, so $r=4$ cm.\par $C=4\times2\times3.14=25.12$ cm.
\dapso{$25.12$ cm.}}
\gap
\begin{probbox}{32}{\stars{1}}
If a circle's radius grows by $2$ cm, by how much does its circumference grow? Does the answer depend on the original radius?
\end{probbox}
\probsol{
The increase is $(r+2)\times2\times3.14-r\times2\times3.14=2\times2\times3.14=12.56$ cm, whatever the original radius.
\dapso{$12.56$ cm; it does not depend on it.}}
\gap
\begin{probbox}{33}{\stars{1}}
A large semicircle has diameter $AB=12$ cm. M is the midpoint of AB, and two small semicircles are drawn on AM and MB. Compare the length of the large arc with the total length of the two small arcs.
\end{probbox}
\probsol{
Large arc: $12\times3.14:2=18.84$ cm.\par Each small arc: $6\times3.14:2=9.42$ cm; together $18.84$ cm.\par They are equal.
\dapso{Equal ($18.84$ cm).}}
\gap
\begin{probbox}{34}{\stars{1}}
Two circles share centre O, with radii $3$ cm and $5$ cm. Find the area between them (the ring).
\end{probbox}
\probsol{
$S=5\times5\times3.14-3\times3\times3.14=(25-9)\times3.14=50.24$ cm$^2$.
\dapso{$50.24$ cm$^2$.}}
\gap
\begin{probbox}{35}{\stars{2}}
\begin{minipage}[t]{0.66\linewidth}
Square ABCD has side $8$ cm. Four quarter circles of radius $4$ cm are drawn with centres at the four corners. Find the shaded area.
\end{minipage}\hfill
\begin{minipage}[t]{0.3\linewidth}\vspace{0pt}\centering
\begin{tikzpicture}[scale=0.24]
\fill[shade] (0,0) rectangle (8,8);
\foreach \x/\y/\a in {0/0/0,8/0/90,8/8/180,0/8/270}{\fill[white] (\x,\y) -- ++(\a:4) arc (\a:\a+90:4) -- cycle; \draw[geo] (\x,\y) ++(\a:4) arc (\a:\a+90:4);}
\draw[geo] (0,0) rectangle (8,8);
\node[lbl,below left] at (0,0){D}; \node[lbl,below right] at (8,0){C}; \node[lbl,above right] at (8,8){B}; \node[lbl,above left] at (0,8){A};
\end{tikzpicture}
\end{minipage}
\end{probbox}
\probsol{
The four quarter circles of radius $4$ cm make one whole circle: $4\times4\times3.14=50.24$ cm$^2$.\par
Shaded part: $8\times8-50.24=13.76$ cm$^2$.
\dapso{$13.76$ cm$^2$.}}
\gap
\begin{probbox}{36}{\stars{2}}
\begin{minipage}[t]{0.66\linewidth}
Square ABCD has side $10$ cm. Two quarter circles of radius $10$ cm are drawn with centres B and D; they overlap in the shaded ``leaf''. Find the area of the leaf.
\end{minipage}\hfill
\begin{minipage}[t]{0.3\linewidth}\vspace{0pt}\centering
\begin{tikzpicture}[scale=0.24]
\fill[shade] (10,0) arc (0:90:10) arc (180:270:10);
\draw[geo] (0,0) rectangle (10,10);
\draw[geo] (10,0) arc (0:90:10); \draw[geo] (0,10) arc (180:270:10);
\draw[dashed,Slate] (0,10)--(10,0);
\node[lbl,below left] at (0,0){D}; \node[lbl,below right] at (10,0){C}; \node[lbl,above right] at (10,10){B}; \node[lbl,above left] at (0,10){A};
\end{tikzpicture}
\end{minipage}
\end{probbox}
\probsol{
Diagonal AC splits the leaf into two equal halves.\par Each half is a quarter circle minus a right-angled triangle:\par
$10\times10\times3.14:4-10\times10:2=78.5-50=28.5$ cm$^2$.\par
Leaf: $28.5\times2=57$ cm$^2$.
\dapso{$57$ cm$^2$.}}
\gap
\begin{probbox}{37}{\stars{1}}
Two circles with the same centre have radii $4$ m and $6$ m. Find the area of the ring between them.
\end{probbox}
\probsol{
Large circle: $6\times6\times3.14=113.04$ m$^2$.\par Small circle: $4\times4\times3.14=50.24$ m$^2$.\par Required area: $113.04-50.24=62.8$ m$^2$.
\dapso{$62.8$ m$^2$.}}
\gap
\begin{probbox}{38}{\stars{1}}
A round flower bed has radius $8$ m, with a path $2$ m wide around it. Find the area of the path.
\end{probbox}
\probsol{
Large circle: $10\times10\times3.14=314$ m$^2$.\par Small circle: $8\times8\times3.14=200.96$ m$^2$.\par Required area: $314-200.96=113.04$ m$^2$.
\dapso{$113.04$ m$^2$.}}
\gap
\begin{probbox}{39}{\stars{1}}
A square has its four corners on a circle of radius $6$ cm. Find the area of the circle outside the square.
\end{probbox}
\probsol{
The square's diagonals are diameters, $12$ cm long and perpendicular.\par It is $4$ right triangles with legs $6$ cm: $6\times6:2\times4=72$ cm$^2$.\par Circle: $6\times6\times3.14=113.04$ cm$^2$.\par Required: $113.04-72=41.04$ cm$^2$.
\dapso{$41.04$ cm$^2$.}}
\gap
\begin{probbox}{40}{\stars{1}}
A small wheel of radius $0.3$ m makes $150$ turns. A large wheel covers the same distance in $90$ turns. Find its radius.
\end{probbox}
\probsol{
Distance: $0.3\times2\times3.14\times150=282.6$ m.\par One large turn: $282.6:90=3.14$ m.\par Radius: $3.14:3.14:2=0.5$ m.
\dapso{$0.5$ m.}}
\gap
\begin{probbox}{41}{\stars{1}}
A circle has the same perimeter as a square of side $15.7$ cm. Which has the larger area, and by how much?
\end{probbox}
\probsol{
Square perimeter: $62.8$ cm.\par Circle radius: $62.8:3.14:2=10$ cm; area $314$ cm$^2$.\par Square area: $15.7\times15.7=246.49$ cm$^2$.\par The circle is larger by $67.51$ cm$^2$.
\dapso{The circle, by $67.51$ cm$^2$.}}
\gap
\begin{probbox}{42}{\stars{1}}
A circle has area $12.56$ cm$^2$. If its radius is doubled, what is the new area?
\end{probbox}
\probsol{
Doubling $r$ multiplies $r\times r$ by $2\times2=4$.\par New area: $12.56\times4=50.24$ cm$^2$.
\dapso{$50.24$ cm$^2$.}}
\gap
\begin{probbox}{43}{\stars{1}}
A circle's radius grows by $1.5$ cm. By how much does its circumference grow?
\end{probbox}
\probsol{
$C=r\times2\times3.14$, so $C$ grows by $1.5\times2\times3.14=9.42$ cm.\par This does not depend on the original radius.
\dapso{$9.42$ cm.}}
\gap
\begin{probbox}{44}{\stars{1}}
An old bicycle has a front wheel of diameter $0.5$ m and a back wheel of diameter $1$ m. When the back wheel turns $100$ times, how many times does the front wheel turn?
\end{probbox}
\probsol{
Distance: $1\times3.14\times100=314$ m.\par One front turn: $0.5\times3.14=1.57$ m.\par Front turns: $314:1.57=200$.
\dapso{$200$ turns.}}
\gap
\begin{probbox}{45}{\stars{1}}
A running track has two straights of $100$ m and two semicircular ends of diameter $60$ m. Find the length of one lap and the area inside.
\end{probbox}
\probsol{
Two semicircles make one circle: $60\times3.14=188.4$ m.\par One lap: $100\times2+188.4=388.4$ m.\par Rectangle: $6\,000$ m$^2$; circle: $30\times30\times3.14=2\,826$ m$^2$.\par Inside area: $8\,826$ m$^2$.
\dapso{$388.4$ m; $8\,826$ m$^2$.}}
\gap
\begin{probbox}{46}{\stars{2}}
A quarter (bounded by two perpendicular radii) is cut from a circular card of radius $10$ cm. Find the area and perimeter of what is left.
\end{probbox}
\probsol{
Circle: $314$ cm$^2$; three quarters: $235.5$ cm$^2$.\par Circumference $62.8$ cm; the remaining arc: $47.1$ cm.\par Perimeter: arc plus two radii, $47.1+20=67.1$ cm.
\dapso{$235.5$ cm$^2$; $67.1$ cm.}}
\gap
\begin{probbox}{47}{\stars{2}}
A goat is tied to an outside corner of a square pen of side $6$ m with a $4$ m rope. It cannot enter the pen. Find the area it can graze.
\end{probbox}
\probsol{
The rope is shorter than the side, so the goat only swings round that corner.\par The pen blocks $\dfrac14$ of the circle, leaving $\dfrac34$ of a circle of radius $4$ m.\par Area: $4\times4\times3.14\times\dfrac34=37.68$ m$^2$.
\dapso{$37.68$ m$^2$.}}
\gap
\begin{probbox}{48}{\stars{1}}
Two circles' circumferences differ by $12.56$ cm, and the larger radius is $3$ times the smaller. Find the area of the larger circle.
\end{probbox}
\probsol{
Difference of radii: $12.56:2:3.14=2$ cm.\par $3-1=2$ parts make $2$ cm, so the radii are $1$ cm and $3$ cm.\par Larger area: $3\times3\times3.14=28.26$ cm$^2$.
\dapso{$28.26$ cm$^2$.}}
\gap
\begin{probbox}{49}{\stars{1}}
A minute hand is $12$ cm long. What area does it sweep in $20$ minutes?
\end{probbox}
\probsol{
$20$ minutes is $\dfrac13$ of a turn.\par Circle: $12\times12\times3.14=452.16$ cm$^2$.\par Swept area: $452.16:3=150.72$ cm$^2$.
\dapso{$150.72$ cm$^2$.}}
\gap
\begin{probbox}{50}{\stars{1}}
A tractor wheel of diameter $0.7$ m makes $120$ turns a minute. How far does the tractor go in $10$ minutes?
\end{probbox}
\probsol{
One turn: $0.7\times3.14=2.198$ m.\par Turns in $10$ minutes: $1\,200$.\par Distance: $2.198\times1\,200=2\,637.6$ m.
\dapso{$2\,637.6$ m.}}
\clearpage
\titlehead
